\documentclass{chsmc}
\chsmcsetup{
	year = 2013,
	part = 1,
	type = A,
	show-answers = false,
}
\begin{document}

\section{填空题：本大题共 8 小题，每小题 8 分，共 64 分.}

\begin{enumerate}
	\item 设集合 $A = \{2,0,1,3\}$，集合
		$B = \{x\mid -x \in A,\, 2 - x^2 \notin A\}$.
		则集合 $B$ 中所有元素的和为 \blankline
	\item 在平面直角坐标系 $xOy$ 中，点 $A$, $B$ 在抛物线 $y^2 = 4x$ 上，
		满足 $\overrightarrow{OA}\cdot\overrightarrow{OB} = -4$，$F$ 是抛物线的焦点.
		则 $S_{\triangle OFA}\cdot S_{\triangle OFB} =$ \blankline
	\item 在 $\triangle ABC$ 中，已知
		$\sin A = 10\sin B\sin C$，$\cos A = 10\cos B\cos C$，
		则 $\tan A$ 的值为 \blankline
	\item 已知正三棱锥 $P$-$ABC$ 底面边长为 $1$，高为 $\sqrt{2}$，
		则其内切球半径为 \blankline
	\item 设 $a$, $b$ 为实数，函数 $f(x) = ax + b$ 满足：
		对任意 $x \in [0,1]$，有 $|f(x)| \leqslant 1$.
		则 $ab$ 的最大值为 \blankline
	\item 从 $1$, $2$, $\cdots$, $20$ 中任取 $5$ 个不同的数，
		其中至少有两个是相邻数的概率为 \blankline
	\item 若实数 $x$, $y$ 满足 $x - 4\sqrt{y} = 2\sqrt{x-y}$，
		则 $x$ 的取值范围是 \blankline
	\item 已知数列 $\{a_n\}$ 共有 $9$ 项，其中 $a_1 = a_9 = 1$，
		且对每个 $i \in \{1,2,\cdots,8\}$，均有
		$\dfrac{a_{i+1}}{a_i} \in \left\{2,1,-\dfrac{1}{2}\right\}$，
		则这样的数列的个数为 \blankline
\end{enumerate}

\section{解答题：本大题共 3 个小题，满分 56 分. 解答应写出文字说明、证明过程或演算步骤.}

\begin{enumerate}[resume]
	\item (本题满分 16 分) 给定正数数列 $\{x_n\}$ 满足
		$S_n \geqslant 2S_{n-1}$，$n = 2,3,\cdots$，这里
		$S_n = x_1 + x_2 + \cdots + x_n$.
		证明：存在常数 $C > 0$，使得 $x_n \geqslant C\cdot 2^n$，$n = 1,2,\cdots$.
	\item (本题满分 20 分) 在平面直角坐标系 $xOy$ 中，椭圆的方程为
		$\dfrac{x^2}{a^2} + \dfrac{y^2}{b^2} = 1$（$a > b > 0$），
		$A_1$, $A_2$ 分别为椭圆的左、右顶点，$F_1$, $F_2$ 分别为椭圆的左、右焦点，
		$P$ 为椭圆上不同于 $A_1$ 和 $A_2$ 的任意一点.
		若平面中两个点 $Q$, $R$ 满足 $QA_1 \perp PA_1$，$QA_2 \perp PA_2$，
		$RF_1 \perp PF_1$，$RF_2 \perp PF_2$，试确定线段 $QR$ 的长度与 $b$ 的大小关系，
		并给出证明.
	\item (本题满分 20 分) 设函数 $f(x) = ax^2 + b$，求所有的正实数对 $(a,b)$，
		使得对任意实数 $x$, $y$，有 $f(xy) + f(x+y) \geqslant f(x)f(y)$.
\end{enumerate}

\end{document}
